
\documentclass[10pt]{article} % For LaTeX2e
%\usepackage{tmlr}
% If accepted, instead use the following line for the camera-ready submission:
%\usepackage[accepted]{tmlr}
% To de-anonymize and remove mentions to TMLR (for example for posting to preprint servers), instead use the following:
\usepackage[preprint]{tmlr}

% Optional math commands from https://github.com/goodfeli/dlbook_notation.
%%%%% NEW MATH DEFINITIONS %%%%%

\usepackage{amsmath,amsfonts,bm}

% Mark sections of captions for referring to divisions of figures
\newcommand{\figleft}{{\em (Left)}}
\newcommand{\figcenter}{{\em (Center)}}
\newcommand{\figright}{{\em (Right)}}
\newcommand{\figtop}{{\em (Top)}}
\newcommand{\figbottom}{{\em (Bottom)}}
\newcommand{\captiona}{{\em (a)}}
\newcommand{\captionb}{{\em (b)}}
\newcommand{\captionc}{{\em (c)}}
\newcommand{\captiond}{{\em (d)}}

% Highlight a newly defined term
\newcommand{\newterm}[1]{{\bf #1}}


% Figure reference, lower-case.
\def\figref#1{figure~\ref{#1}}
% Figure reference, capital. For start of sentence
\def\Figref#1{Figure~\ref{#1}}
\def\twofigref#1#2{figures \ref{#1} and \ref{#2}}
\def\quadfigref#1#2#3#4{figures \ref{#1}, \ref{#2}, \ref{#3} and \ref{#4}}
% Section reference, lower-case.
\def\secref#1{section~\ref{#1}}
% Section reference, capital.
\def\Secref#1{Section~\ref{#1}}
% Reference to two sections.
\def\twosecrefs#1#2{sections \ref{#1} and \ref{#2}}
% Reference to three sections.
\def\secrefs#1#2#3{sections \ref{#1}, \ref{#2} and \ref{#3}}
% Reference to an equation, lower-case.
\def\eqref#1{equation~\ref{#1}}
% Reference to an equation, upper case
\def\Eqref#1{Equation~\ref{#1}}
% A raw reference to an equation---avoid using if possible
\def\plaineqref#1{\ref{#1}}
% Reference to a chapter, lower-case.
\def\chapref#1{chapter~\ref{#1}}
% Reference to an equation, upper case.
\def\Chapref#1{Chapter~\ref{#1}}
% Reference to a range of chapters
\def\rangechapref#1#2{chapters\ref{#1}--\ref{#2}}
% Reference to an algorithm, lower-case.
\def\algref#1{algorithm~\ref{#1}}
% Reference to an algorithm, upper case.
\def\Algref#1{Algorithm~\ref{#1}}
\def\twoalgref#1#2{algorithms \ref{#1} and \ref{#2}}
\def\Twoalgref#1#2{Algorithms \ref{#1} and \ref{#2}}
% Reference to a part, lower case
\def\partref#1{part~\ref{#1}}
% Reference to a part, upper case
\def\Partref#1{Part~\ref{#1}}
\def\twopartref#1#2{parts \ref{#1} and \ref{#2}}

\def\ceil#1{\lceil #1 \rceil}
\def\floor#1{\lfloor #1 \rfloor}
\def\1{\bm{1}}
\newcommand{\train}{\mathcal{D}}
\newcommand{\valid}{\mathcal{D_{\mathrm{valid}}}}
\newcommand{\test}{\mathcal{D_{\mathrm{test}}}}

\def\eps{{\epsilon}}


% Random variables
\def\reta{{\textnormal{$\eta$}}}
\def\ra{{\textnormal{a}}}
\def\rb{{\textnormal{b}}}
\def\rc{{\textnormal{c}}}
\def\rd{{\textnormal{d}}}
\def\re{{\textnormal{e}}}
\def\rf{{\textnormal{f}}}
\def\rg{{\textnormal{g}}}
\def\rh{{\textnormal{h}}}
\def\ri{{\textnormal{i}}}
\def\rj{{\textnormal{j}}}
\def\rk{{\textnormal{k}}}
\def\rl{{\textnormal{l}}}
% rm is already a command, just don't name any random variables m
\def\rn{{\textnormal{n}}}
\def\ro{{\textnormal{o}}}
\def\rp{{\textnormal{p}}}
\def\rq{{\textnormal{q}}}
\def\rr{{\textnormal{r}}}
\def\rs{{\textnormal{s}}}
\def\rt{{\textnormal{t}}}
\def\ru{{\textnormal{u}}}
\def\rv{{\textnormal{v}}}
\def\rw{{\textnormal{w}}}
\def\rx{{\textnormal{x}}}
\def\ry{{\textnormal{y}}}
\def\rz{{\textnormal{z}}}

% Random vectors
\def\rvepsilon{{\mathbf{\epsilon}}}
\def\rvtheta{{\mathbf{\theta}}}
\def\rva{{\mathbf{a}}}
\def\rvb{{\mathbf{b}}}
\def\rvc{{\mathbf{c}}}
\def\rvd{{\mathbf{d}}}
\def\rve{{\mathbf{e}}}
\def\rvf{{\mathbf{f}}}
\def\rvg{{\mathbf{g}}}
\def\rvh{{\mathbf{h}}}
\def\rvu{{\mathbf{i}}}
\def\rvj{{\mathbf{j}}}
\def\rvk{{\mathbf{k}}}
\def\rvl{{\mathbf{l}}}
\def\rvm{{\mathbf{m}}}
\def\rvn{{\mathbf{n}}}
\def\rvo{{\mathbf{o}}}
\def\rvp{{\mathbf{p}}}
\def\rvq{{\mathbf{q}}}
\def\rvr{{\mathbf{r}}}
\def\rvs{{\mathbf{s}}}
\def\rvt{{\mathbf{t}}}
\def\rvu{{\mathbf{u}}}
\def\rvv{{\mathbf{v}}}
\def\rvw{{\mathbf{w}}}
\def\rvx{{\mathbf{x}}}
\def\rvy{{\mathbf{y}}}
\def\rvz{{\mathbf{z}}}

% Elements of random vectors
\def\erva{{\textnormal{a}}}
\def\ervb{{\textnormal{b}}}
\def\ervc{{\textnormal{c}}}
\def\ervd{{\textnormal{d}}}
\def\erve{{\textnormal{e}}}
\def\ervf{{\textnormal{f}}}
\def\ervg{{\textnormal{g}}}
\def\ervh{{\textnormal{h}}}
\def\ervi{{\textnormal{i}}}
\def\ervj{{\textnormal{j}}}
\def\ervk{{\textnormal{k}}}
\def\ervl{{\textnormal{l}}}
\def\ervm{{\textnormal{m}}}
\def\ervn{{\textnormal{n}}}
\def\ervo{{\textnormal{o}}}
\def\ervp{{\textnormal{p}}}
\def\ervq{{\textnormal{q}}}
\def\ervr{{\textnormal{r}}}
\def\ervs{{\textnormal{s}}}
\def\ervt{{\textnormal{t}}}
\def\ervu{{\textnormal{u}}}
\def\ervv{{\textnormal{v}}}
\def\ervw{{\textnormal{w}}}
\def\ervx{{\textnormal{x}}}
\def\ervy{{\textnormal{y}}}
\def\ervz{{\textnormal{z}}}

% Random matrices
\def\rmA{{\mathbf{A}}}
\def\rmB{{\mathbf{B}}}
\def\rmC{{\mathbf{C}}}
\def\rmD{{\mathbf{D}}}
\def\rmE{{\mathbf{E}}}
\def\rmF{{\mathbf{F}}}
\def\rmG{{\mathbf{G}}}
\def\rmH{{\mathbf{H}}}
\def\rmI{{\mathbf{I}}}
\def\rmJ{{\mathbf{J}}}
\def\rmK{{\mathbf{K}}}
\def\rmL{{\mathbf{L}}}
\def\rmM{{\mathbf{M}}}
\def\rmN{{\mathbf{N}}}
\def\rmO{{\mathbf{O}}}
\def\rmP{{\mathbf{P}}}
\def\rmQ{{\mathbf{Q}}}
\def\rmR{{\mathbf{R}}}
\def\rmS{{\mathbf{S}}}
\def\rmT{{\mathbf{T}}}
\def\rmU{{\mathbf{U}}}
\def\rmV{{\mathbf{V}}}
\def\rmW{{\mathbf{W}}}
\def\rmX{{\mathbf{X}}}
\def\rmY{{\mathbf{Y}}}
\def\rmZ{{\mathbf{Z}}}

% Elements of random matrices
\def\ermA{{\textnormal{A}}}
\def\ermB{{\textnormal{B}}}
\def\ermC{{\textnormal{C}}}
\def\ermD{{\textnormal{D}}}
\def\ermE{{\textnormal{E}}}
\def\ermF{{\textnormal{F}}}
\def\ermG{{\textnormal{G}}}
\def\ermH{{\textnormal{H}}}
\def\ermI{{\textnormal{I}}}
\def\ermJ{{\textnormal{J}}}
\def\ermK{{\textnormal{K}}}
\def\ermL{{\textnormal{L}}}
\def\ermM{{\textnormal{M}}}
\def\ermN{{\textnormal{N}}}
\def\ermO{{\textnormal{O}}}
\def\ermP{{\textnormal{P}}}
\def\ermQ{{\textnormal{Q}}}
\def\ermR{{\textnormal{R}}}
\def\ermS{{\textnormal{S}}}
\def\ermT{{\textnormal{T}}}
\def\ermU{{\textnormal{U}}}
\def\ermV{{\textnormal{V}}}
\def\ermW{{\textnormal{W}}}
\def\ermX{{\textnormal{X}}}
\def\ermY{{\textnormal{Y}}}
\def\ermZ{{\textnormal{Z}}}

% Vectors
\def\vzero{{\bm{0}}}
\def\vone{{\bm{1}}}
\def\vmu{{\bm{\mu}}}
\def\vtheta{{\bm{\theta}}}
\def\va{{\bm{a}}}
\def\vb{{\bm{b}}}
\def\vc{{\bm{c}}}
\def\vd{{\bm{d}}}
\def\ve{{\bm{e}}}
\def\vf{{\bm{f}}}
\def\vg{{\bm{g}}}
\def\vh{{\bm{h}}}
\def\vi{{\bm{i}}}
\def\vj{{\bm{j}}}
\def\vk{{\bm{k}}}
\def\vl{{\bm{l}}}
\def\vm{{\bm{m}}}
\def\vn{{\bm{n}}}
\def\vo{{\bm{o}}}
\def\vp{{\bm{p}}}
\def\vq{{\bm{q}}}
\def\vr{{\bm{r}}}
\def\vs{{\bm{s}}}
\def\vt{{\bm{t}}}
\def\vu{{\bm{u}}}
\def\vv{{\bm{v}}}
\def\vw{{\bm{w}}}
\def\vx{{\bm{x}}}
\def\vy{{\bm{y}}}
\def\vz{{\bm{z}}}

% Elements of vectors
\def\evalpha{{\alpha}}
\def\evbeta{{\beta}}
\def\evepsilon{{\epsilon}}
\def\evlambda{{\lambda}}
\def\evomega{{\omega}}
\def\evmu{{\mu}}
\def\evpsi{{\psi}}
\def\evsigma{{\sigma}}
\def\evtheta{{\theta}}
\def\eva{{a}}
\def\evb{{b}}
\def\evc{{c}}
\def\evd{{d}}
\def\eve{{e}}
\def\evf{{f}}
\def\evg{{g}}
\def\evh{{h}}
\def\evi{{i}}
\def\evj{{j}}
\def\evk{{k}}
\def\evl{{l}}
\def\evm{{m}}
\def\evn{{n}}
\def\evo{{o}}
\def\evp{{p}}
\def\evq{{q}}
\def\evr{{r}}
\def\evs{{s}}
\def\evt{{t}}
\def\evu{{u}}
\def\evv{{v}}
\def\evw{{w}}
\def\evx{{x}}
\def\evy{{y}}
\def\evz{{z}}

% Matrix
\def\mA{{\bm{A}}}
\def\mB{{\bm{B}}}
\def\mC{{\bm{C}}}
\def\mD{{\bm{D}}}
\def\mE{{\bm{E}}}
\def\mF{{\bm{F}}}
\def\mG{{\bm{G}}}
\def\mH{{\bm{H}}}
\def\mI{{\bm{I}}}
\def\mJ{{\bm{J}}}
\def\mK{{\bm{K}}}
\def\mL{{\bm{L}}}
\def\mM{{\bm{M}}}
\def\mN{{\bm{N}}}
\def\mO{{\bm{O}}}
\def\mP{{\bm{P}}}
\def\mQ{{\bm{Q}}}
\def\mR{{\bm{R}}}
\def\mS{{\bm{S}}}
\def\mT{{\bm{T}}}
\def\mU{{\bm{U}}}
\def\mV{{\bm{V}}}
\def\mW{{\bm{W}}}
\def\mX{{\bm{X}}}
\def\mY{{\bm{Y}}}
\def\mZ{{\bm{Z}}}
\def\mBeta{{\bm{\beta}}}
\def\mPhi{{\bm{\Phi}}}
\def\mLambda{{\bm{\Lambda}}}
\def\mSigma{{\bm{\Sigma}}}

% Tensor
\DeclareMathAlphabet{\mathsfit}{\encodingdefault}{\sfdefault}{m}{sl}
\SetMathAlphabet{\mathsfit}{bold}{\encodingdefault}{\sfdefault}{bx}{n}
\newcommand{\tens}[1]{\bm{\mathsfit{#1}}}
\def\tA{{\tens{A}}}
\def\tB{{\tens{B}}}
\def\tC{{\tens{C}}}
\def\tD{{\tens{D}}}
\def\tE{{\tens{E}}}
\def\tF{{\tens{F}}}
\def\tG{{\tens{G}}}
\def\tH{{\tens{H}}}
\def\tI{{\tens{I}}}
\def\tJ{{\tens{J}}}
\def\tK{{\tens{K}}}
\def\tL{{\tens{L}}}
\def\tM{{\tens{M}}}
\def\tN{{\tens{N}}}
\def\tO{{\tens{O}}}
\def\tP{{\tens{P}}}
\def\tQ{{\tens{Q}}}
\def\tR{{\tens{R}}}
\def\tS{{\tens{S}}}
\def\tT{{\tens{T}}}
\def\tU{{\tens{U}}}
\def\tV{{\tens{V}}}
\def\tW{{\tens{W}}}
\def\tX{{\tens{X}}}
\def\tY{{\tens{Y}}}
\def\tZ{{\tens{Z}}}


% Graph
\def\gA{{\mathcal{A}}}
\def\gB{{\mathcal{B}}}
\def\gC{{\mathcal{C}}}
\def\gD{{\mathcal{D}}}
\def\gE{{\mathcal{E}}}
\def\gF{{\mathcal{F}}}
\def\gG{{\mathcal{G}}}
\def\gH{{\mathcal{H}}}
\def\gI{{\mathcal{I}}}
\def\gJ{{\mathcal{J}}}
\def\gK{{\mathcal{K}}}
\def\gL{{\mathcal{L}}}
\def\gM{{\mathcal{M}}}
\def\gN{{\mathcal{N}}}
\def\gO{{\mathcal{O}}}
\def\gP{{\mathcal{P}}}
\def\gQ{{\mathcal{Q}}}
\def\gR{{\mathcal{R}}}
\def\gS{{\mathcal{S}}}
\def\gT{{\mathcal{T}}}
\def\gU{{\mathcal{U}}}
\def\gV{{\mathcal{V}}}
\def\gW{{\mathcal{W}}}
\def\gX{{\mathcal{X}}}
\def\gY{{\mathcal{Y}}}
\def\gZ{{\mathcal{Z}}}

% Sets
\def\sA{{\mathbb{A}}}
\def\sB{{\mathbb{B}}}
\def\sC{{\mathbb{C}}}
\def\sD{{\mathbb{D}}}
% Don't use a set called E, because this would be the same as our symbol
% for expectation.
\def\sF{{\mathbb{F}}}
\def\sG{{\mathbb{G}}}
\def\sH{{\mathbb{H}}}
\def\sI{{\mathbb{I}}}
\def\sJ{{\mathbb{J}}}
\def\sK{{\mathbb{K}}}
\def\sL{{\mathbb{L}}}
\def\sM{{\mathbb{M}}}
\def\sN{{\mathbb{N}}}
\def\sO{{\mathbb{O}}}
\def\sP{{\mathbb{P}}}
\def\sQ{{\mathbb{Q}}}
\def\sR{{\mathbb{R}}}
\def\sS{{\mathbb{S}}}
\def\sT{{\mathbb{T}}}
\def\sU{{\mathbb{U}}}
\def\sV{{\mathbb{V}}}
\def\sW{{\mathbb{W}}}
\def\sX{{\mathbb{X}}}
\def\sY{{\mathbb{Y}}}
\def\sZ{{\mathbb{Z}}}

% Entries of a matrix
\def\emLambda{{\Lambda}}
\def\emA{{A}}
\def\emB{{B}}
\def\emC{{C}}
\def\emD{{D}}
\def\emE{{E}}
\def\emF{{F}}
\def\emG{{G}}
\def\emH{{H}}
\def\emI{{I}}
\def\emJ{{J}}
\def\emK{{K}}
\def\emL{{L}}
\def\emM{{M}}
\def\emN{{N}}
\def\emO{{O}}
\def\emP{{P}}
\def\emQ{{Q}}
\def\emR{{R}}
\def\emS{{S}}
\def\emT{{T}}
\def\emU{{U}}
\def\emV{{V}}
\def\emW{{W}}
\def\emX{{X}}
\def\emY{{Y}}
\def\emZ{{Z}}
\def\emSigma{{\Sigma}}

% entries of a tensor
% Same font as tensor, without \bm wrapper
\newcommand{\etens}[1]{\mathsfit{#1}}
\def\etLambda{{\etens{\Lambda}}}
\def\etA{{\etens{A}}}
\def\etB{{\etens{B}}}
\def\etC{{\etens{C}}}
\def\etD{{\etens{D}}}
\def\etE{{\etens{E}}}
\def\etF{{\etens{F}}}
\def\etG{{\etens{G}}}
\def\etH{{\etens{H}}}
\def\etI{{\etens{I}}}
\def\etJ{{\etens{J}}}
\def\etK{{\etens{K}}}
\def\etL{{\etens{L}}}
\def\etM{{\etens{M}}}
\def\etN{{\etens{N}}}
\def\etO{{\etens{O}}}
\def\etP{{\etens{P}}}
\def\etQ{{\etens{Q}}}
\def\etR{{\etens{R}}}
\def\etS{{\etens{S}}}
\def\etT{{\etens{T}}}
\def\etU{{\etens{U}}}
\def\etV{{\etens{V}}}
\def\etW{{\etens{W}}}
\def\etX{{\etens{X}}}
\def\etY{{\etens{Y}}}
\def\etZ{{\etens{Z}}}

% The true underlying data generating distribution
\newcommand{\pdata}{p_{\rm{data}}}
% The empirical distribution defined by the training set
\newcommand{\ptrain}{\hat{p}_{\rm{data}}}
\newcommand{\Ptrain}{\hat{P}_{\rm{data}}}
% The model distribution
\newcommand{\pmodel}{p_{\rm{model}}}
\newcommand{\Pmodel}{P_{\rm{model}}}
\newcommand{\ptildemodel}{\tilde{p}_{\rm{model}}}
% Stochastic autoencoder distributions
\newcommand{\pencode}{p_{\rm{encoder}}}
\newcommand{\pdecode}{p_{\rm{decoder}}}
\newcommand{\precons}{p_{\rm{reconstruct}}}

\newcommand{\laplace}{\mathrm{Laplace}} % Laplace distribution

\newcommand{\E}{\mathbb{E}}
\newcommand{\Ls}{\mathcal{L}}
\newcommand{\R}{\mathbb{R}}
\newcommand{\emp}{\tilde{p}}
\newcommand{\lr}{\alpha}
\newcommand{\reg}{\lambda}
\newcommand{\rect}{\mathrm{rectifier}}
\newcommand{\softmax}{\mathrm{softmax}}
\newcommand{\sigmoid}{\sigma}
\newcommand{\softplus}{\zeta}
\newcommand{\KL}{D_{\mathrm{KL}}}
\newcommand{\Var}{\mathrm{Var}}
\newcommand{\standarderror}{\mathrm{SE}}
\newcommand{\Cov}{\mathrm{Cov}}
% Wolfram Mathworld says $L^2$ is for function spaces and $\ell^2$ is for vectors
% But then they seem to use $L^2$ for vectors throughout the site, and so does
% wikipedia.
\newcommand{\normlzero}{L^0}
\newcommand{\normlone}{L^1}
\newcommand{\normltwo}{L^2}
\newcommand{\normlp}{L^p}
\newcommand{\normmax}{L^\infty}

\newcommand{\parents}{Pa} % See usage in notation.tex. Chosen to match Daphne's book.

\DeclareMathOperator*{\argmax}{arg\,max}
\DeclareMathOperator*{\argmin}{arg\,min}

\DeclareMathOperator{\sign}{sign}
\DeclareMathOperator{\Tr}{Tr}
\let\ab\allowbreak


\usepackage{hyperref}
\usepackage{url}


\title{Sull'affidabilità delle reti neurali}
% [Eng] A formal methodology for the Reliability of Neural Networks
% Authors must not appear in the submitted version. They should be hidden
% as long as the tmlr package is used without the [accepted] or [preprint] options.
% Non-anonymous submissions will be rejected without review.

\author{\name Rosario Pio Gnazzo \email rosariopiognazzo@gmail.com \\
      \addr Department of Computer Science\\
      University of Salerno}

% The \author macro works with any number of authors. Use \AND 
% to separate the names and addresses of multiple authors.

\newcommand{\fix}{\marginpar{FIX}}
\newcommand{\new}{\marginpar{NEW}}

\def\month{MM}  % Insert correct month for camera-ready version
\def\year{YYYY} % Insert correct year for camera-ready version
\def\openreview{\url{https://openreview.net/forum?id=XXXX}} % Insert correct link to OpenReview for camera-ready version


\begin{document}


\maketitle

\begin{abstract}
Sebbene il Deep Learning abbia raggiunto prestazioni eccellenti in numerosi ambiti, manca di \textit{metriche formali} per quantificare il margine di sicurezza e la robustezza previsionale in contesti \textit{safety-critical}. Normalmente ci si focalizzata sul valutare un modello sulle sole metriche di valutazione (performance-based), senza tener conto di una visione sull'\textit{affidabilità} intrinseca del sistema.

Questo lavoro prova a studiare un metodo di analisi matematico-statistico integrando la \textit{Teoria dell'Affidabilità classica} nel \textit{Deep Learning}. L'idea è, modellando i layer neurali come sistemi stocastici di tipo \textit{k-out-of-n}, di focalizzarsi sulle \textit{attivazioni neuronali} piuttosto che sui pesi, utilizzando la distribuzione Beta-Binomiale per rappresentare il comportamento (inteso come la probabilità di essere \textit{attivi} quando la rete ottiene un successo) dei neuroni.

Il metodo permette di stimare \textit{parametri di affidabilità puntuali}, \textit{soglie di attivazione} minime $s_{(k)}$ e \textit{metriche operative} $\rho_{(k)}$. Inoltre, per gestire l'incertezza delle stime ed avere un'analisi più robusta, vengono derivati in modo analitico gli \textit{intervalli di confidenza}.
\end{abstract}

\section{Introduzione}
\label{intro}
Uno dei problemi che si riscontra nel deep learning è che non si riesce a stimare una robusta previsione di come sarà il comportamento di un modello sui dati di test. Normalmente ci si concentra sul valutare l'\textit{accuratezza} del modello (oppure altre metriche), senza avere contezza di un \textit{margine di sicurezza} della rete, ossia: \textit{quanto sia affidabile nel svolgere il task}. In questo studio andrò a presentare una modellazione che propone esattamente questo, sfruttando le idee già note nel campo dell'affidabilità dei sistemi multicomponente\footnote{La teoria dell'affidabilità è da anni studiata in ingegneria ed informatica come strumento per valutare un sistema sulla base delle proprietà delle componenti di cui è formato. L'intuizione da cui parte il lavoro è quella di considerare una rete neurale come un sistema di componenti che collegati fra loro hanno delle proprietà statistiche che descrivono l'affidabilità delle rete.}. 
L'approccio si focalizza sulle attivazioni neuronali, trattando i singoli neuroni come se fossero componenti fisiche di un sistema.

\section{Definizioni}
\label{assunzioni}
Prima di focalizzarci sull'approccio definiamo gli oggetti che lo compongono e le loro caratteristiche.

\subsection{Neural Network}
\label{A1-FeedForward}
Definiamo con la notazione $f: X^{n} \xrightarrow{\mathbb{H}^{L}} Y^{m}$ una rete neurale tale che:
\begin{itemize}
    \item $X^{n} = \{x_{1}, x_{2}, \dots, x_{n}\}$, identifica \textit{layer di input} composto da $n$ neuroni;
    \item $\mathbb{H}^{L} = \{H^{(1)}_{n_{1}}, H^{(2)}_{n_{2}}, \dots, H^{(L)}_{n_{L}}\}$ che identifica l'\textit{insieme degli} $L$ \textit{layers interni} composti ciascuno da $n_{l}$ neuroni;
    \begin{equation}
        H^{(l)}_{n_{l}} = \{h^{(l)}_{1}, h^{(l)}_{2}, \dots, h^{(l)}_{n_{l}}\}\qquad l = 1, \dots, L
    \end{equation}
    \item $Y^{m} = \{y_{1}, y_{2}, \dots, y_{m}\}$, identifica \textit{layer di output} composto da $m$ neuroni.
\end{itemize}
Data questa formalizzazione definiamo
\begin{itemize}
    \item $z^{(k)}_{i}$ l'\textit{Attivazione del neurone} \textit{i-esimo} che fa parte del \textit{k-esimo} layer post-utilizzo della funzione di attivazione $\sigma$;
\end{itemize}

\subsection{Attivazione neuronale}
\label{A2-Attivazione}
Per comprendere la definizione di affidabilità che sarà data più avanti nel lavoro, è necessario definire cosa rappresenta in termini probabilistici lo stato di \textit{attivazione} per un neurone. Con $A^{(k)}_{i}$ indichiamo che il neurone \textit{i} del \textit{k-esimo} layer è attivo. Questa definizione può essere formalizzata come una funzione binaria
\begin{equation}
    A^{(k)}_{i} = \mathbb{1}[z^{(k)}_{i} > \tau^{(l)}] = 
    \begin{cases} 
        1, & \mbox{se }z^{(k)}_{i} > \tau^{(l)} \\ 
        0, & \mbox{altrimenti}
    \end{cases}
\end{equation}
dove con $\tau^{(l)}$ si indica la soglia di attivazione che dipende dal layer.\footnote{e.g. Se $\tau^{(l)} = 0$ allora la funzione di attivazione che stiamo considerando è \textit{ReLU}}.

\subsection{Componente del sistema}
\label{A3-Component}
Dalla definizione [\ref{A2-Attivazione}] possiamo descrivere un neurone come componente di un sistema osservato sotto la propria condizione di attivazione. Dato un neurone $n^{(k)}_{i}$ appartenente a qualsiasi layer \textit{k} di una rete neurale [\ref{A1-FeedForward}], definiamo l'evento $B = \{\text{Osservare output \textit{corretto} dalla rete}\}$, allora la componente $C^{(k)}_{i}$ relativa al neurone $n^{(k)}_{i}$ sarà descritta da una \textit{v.a. Bernoulli}
\begin{equation} 
    C^{(k)}_{i} \sim Bernoulli(R^{(k)}_{i})
\end{equation}
dove $R^{(k)}_{i}$\footnote{Empiricamente si calcola usando la definizione frequentista di probabilità $\frac{\mbox{Num. Attivazioni con Output Corretto}}{\mbox{Num. Tot. Output Corretto}}$}
misura la \textit{probabilità di partecipazione al successo}, ossia la probabilità che il neurone \textit{i} si accenda quando la rete produce un risultato corretto.
\begin{equation}
    R^{(k)}_{i} = \mathbb{P}\{A^{(k)}_{i} = 1\ |\ B\}
\end{equation}

\section{Metodologia}
\label{A5-k-outof-}
In questa sezione andrò a descrivere la modellazione statistica alla base dell'approccio, per poi ricavarci le metriche di valutazione dell'affidabilità.  

\subsection{Affidabilità del Layer}
Partendo dall'assunto che, una rete neurale in funzione la possiamo assimilare all'insieme delle attivazioni dei suoi neuroni, allora possiamo affermare che \textit{l'affidabilità di un layer è funzione del numero di attivazioni dei suoi neuroni}. Più nello specifico è funzione della affidabilità di un \textit{sotto-gruppo} di questi nell'essere attivi quando avviene il successo.
Nella teoria dell'affidabilità, questa caratteristica è rappresentabile tramite un sistema \textit{k-out-of-n}
\begin{equation}
    \mathcal{R}_{sistema} = \sum_{j=k}^{n}\binom{n}{j} R^{j}(1-R)^{(n-j)} \qquad k=1,\dots, n
\end{equation}
dove 
\begin{align}
    &\binom{n}{k} R^{k} \longrightarrow \text{Probabilità delle $k$ componenti su $n$ di funzionare correttamente} \\
    &\binom{n}{k} (1-R)^{(n-k)} \longrightarrow \text{Probabilità delle restanti $(n-k)$ componenti su $n$ di non funzionare}
\end{align}
Questa formulazione utilizza la \textit{v.a. Binomiale} per modellare la probabilità delle componenti. Tuttavia, per l'affidabilità di un layer di una rete neurale, non è possibile utilizzare la Binomiale in quanto impone le condizioni \textit{i.i.d.} alle v.a. che descrive. Se, semplificando il concetto di rete neurale, la condizione di \textit{indipendenza} potremmo accettarla; la condizione di \textit{identicamente distribuite} non può essere accettata in quanto imporrebbe ad ogni componente del layer (neurone) di avere la \textit{stessa} distribuzione di affidabilità.
Per questo motivo, la soluzione che ho deciso di utilizzare è di usare la versione generalizzata della Binomiale, ossia la \textit{v.a. Beta-Binomiale}.\footnote{In teoria delle probabilità la distribuzione Beta-Binomiale è una famiglia di distribuzioni di probabilità discrete che può essere vista come generalizzazione della distribuzione binomiale e della distribuzione Beta. Descrive la distribuzione del numero di successi su \textit{n} esperimenti indipendenti di tipo sì/no, ma, contrariamente alla distribuzione Binomiale, la probabilità di successo non è un parametro noto, ma è un valore incerto distribuito come una variabile casuale Beta.}

Se $L_{(k)}\sim BetaBin(n_{(k)},\alpha_{(k)},\beta_{(k)})$ è un layer-$k$ definito come una variabile casuale che si distribuisce come una v.a. Beta-Binomiale con i parametri $n_{(k)},\alpha_{(k)},\beta_{(k)}$. Allora con $x \ge 0$ che indica il \textit{numero di componenti attivi} si ha
\begin{equation}\label{BetaBin}
    \mathbb{P}\{L_{(k)} = x\} = \binom{n_{(k)}}{x}\dfrac{\textbf{B}(\alpha_{(k)}+x\ ,\ \beta_{(k)}+n_{(k)}-x)}{\textbf{B}(\alpha_{(k)}\ ,\ \beta_{(k)})}
\end{equation}
con
\begin{equation}
    \mathbb{E}[X] = n\dfrac{\alpha}{\alpha+\beta} \qquad 
    \mathbb{V}[X] = n\dfrac{\alpha[n(1+\alpha)+\beta]}{(\alpha+\beta)(1+\alpha+\beta)}
\end{equation}
dove \textbf{B} è la funzione di \textit{Beta di Eulero}:
$\textbf{B}(a, b) = \int_{0}^{1} y^{\ a-1}(1-y)^{\ b-1} dy$ \footnote{Questa distribuzione governa la probabilità 
\textit{p} di un processo di Bernoulli a posteriori dell'osservazione di $(\alpha-1)$ "successi" e $(\beta-1)$ "fallimenti", quando $p$ è a priori distribuita uniformemente tra $[0,1]$}.

La Beta-Binomiale gestisce la sovra-dispersione ossia l'\textit{alta variabilità}, infatti se in una Binomiale classica la varianza è vincolata alla media, in un layer è molto facile avere neuroni che si attivano sempre e altri invece mai e ciò comporta una varianza alta. La Beta-Binomiale riesce a catturare proprio questa \textit{eterogeneità} dei neuroni all'interno dello stesso layer.

A questo punto possiamo definire l'\textit{affidabilità del layer} come nel modello \textit{k-out-of-n} sostituendo la distribuzione Binomiale con la Beta-Binomiale
\begin{align}\label{Affidabilità Layer}
    \mathcal{R}_{(k)}(n_{(k)},\alpha_{(k)}, \beta_{(k)}) &= \sum_{j=s_{(k)}}^{n_{(k)}} \binom{n_{(k)}}{j}\dfrac{\textbf{B}(\alpha_{(k)}+j\ ,\ \beta_{(k)}+n_{(k)}-j)}{\textbf{B}(\alpha_{(k)}\ ,\ \beta_{(k)})}\\
    &=\sum_{j=s_{(k)}}^{n_{(k)}} \binom{n_{(k)}}{j}\dfrac{\int_{0}^{1} y^{\ (\alpha_{(k)}+j-1)}(1-y)^{\ (\beta_{(k)}+n_{(k)}-j-1)} dy}{\int_{0}^{1} y^{\ \alpha_{(k)}-1}(1-y)^{\ \beta_{(k)}-1} dy}
\end{align}
Affidabilità di un layer della rete è definita come la \textit{massa di probabilità totale} di successo, ossia la probabilità di osservare un comportamento: \textit{l'attivazione di un sotto-gruppo di nodi pari ad almeno} $s_{(k)}$, che sia statisticamente \textit{significativo} quando la rete restituisce un output corretto.

\subsection{Stima dei parametri $\alpha$ e $\beta$}
La distribuzione Beta-Binomiale [\ref{BetaBin}] usata per descrivere un layer, dipende da tre parametri: $\alpha_{(k)}, \beta_{(k)}$ e $n_{(k)}$. Per poter derivare una stima della distribuzione di probabilità del layer è necessario dunque stimare tali parametri. Più nello specifico, nel nostro caso dobbiamo ottenere una stima dei parametri $\alpha_{(k)}$ e $\beta_{(k)}$, in quanto $n_{(k)}$ essendo il numero di nodi presenti nel layer è un valore noto nel momento che si costruisce la rete.

Per stimare i parametri ho utilizzato il \textit{Metodo dei Momenti} in quanto è spesso preferibile alla massima verosimiglianza \footnote{Le stime di massima verosimiglianza in forma chiusa sono poco pratiche, dato che la pdf è costituita da funzioni comuni (funzione gamma e/o funzioni beta), esse possono essere facilmente ricavate tramite ottimizzazione numerica diretta. Le stime di massima verosimiglianza da dati empirici possono essere calcolate utilizzando metodi generali per l'adattamento di distribuzioni di Pólya multinomiali.} per stabilità computazionale e interpretabilità sopratutto per distribuzioni più complesse. Il metodo dei momenti si basa sull'uguaglianza tra l'equazione teoriche dei \textit{momenti della popolazione} ($\mu_{r}$) con quelle \textit{campionarie} ($m_{r}$). Nel nostro caso, siccome i parametri da stimare sono due, costruiremo un sistema di due equazioni composto dall'uguaglianza tra la media campionaria $m_{1}$ e il valore atteso della popolazione $\mu_{1}$ ed la varianza campionaria $m_{2}$ e la varianza della popolazione $\mu_{2}$.
\begin{equation}
    \begin{cases}
        m_{1} &= \mu_{1}\\
        m_{2} &= \mu_{2}
    \end{cases} \qquad
    \begin{cases}
        \bar{\mathbf{x}}_{n} &= \mathbb{E}[L_{(k)}]\\
        S^{2} &= \mathbb{V}[L_{(k)}]
    \end{cases} \qquad
    \begin{cases}
        \bar{\mathbf{x}}_{n} &= n_{(k)}\dfrac{\alpha_{(k)}}{\alpha_{(k)}+\beta_{(k)}}\\
        S^{2} &= n_{(k)}\dfrac{\alpha_{(k)}[n_{(k)}(1+\alpha_{(k)})+\beta_{(k)}]}{(\alpha_{(k)}+\beta_{(k)})(1+\alpha_{(k)}+\beta_{(k)})}
    \end{cases}
\end{equation}
La prima equazione si esplicita per $\alpha_{(k)}$ e si ottiene $\frac{\bar{\mathbf{x}}_{n}\beta_{(k)}}{n_{(k)}-\bar{\mathbf{x}}_{n}}$, sostituendo questo valore nella seconda equazione ed invertendola per $\beta_{(k)}$ otteniamo le \textit{stime puntuali} dei parametri
\begin{equation}\label{stime parametri}
    \begin{cases}
        \hat\alpha_{(k)} &= \dfrac{n_{(k)}m_{1}-m_{2}}{n_{(k)}(\frac{m_{2}}{m_{1}}-{m_{1}} -1)+m_{1}}\\ \\
        \hat\beta_{(k)} &= \dfrac{(n_{(k)}-m_{1})(n_{(k)}-\frac{m_{2}}{m_{1}})}{n_{(k)}(\frac{m_{2}}{m_{1}}-m_{1} -1)+m_{1}}
    \end{cases} \quad \mbox{con } m_{1} = \bar{\mathbf{x}}_{n}\mbox{ e } m_{2} = S^{2}
\end{equation}
Ora possiamo utilizzarle per stimare la distribuzione di probabilità della Beta-Binomiale ed utilizzarla per studiare la soglia $s_{(k)}$.\footnote{Nelle sezioni successive del lavoro verrà ampliata l'analisi gestendo l'incertezza delle stime costruendo gli intervalli di confidenza.}

\subsection{Derivazione soglia $s_{(k)}$}
Dalla definizione dell'affidabilità del layer [\ref{Affidabilità Layer}] abbiamo detto che il parametro $s_{(k)}$ rappresenta il più piccolo numero di neuroni che devono essere attivi affinché l'attivazione generale del layer sia statisticamente significativa rispetto al comportamento di successo della rete. In termini formali diremo che $s_{(k)}$ \textit{è il più piccolo valore intero per cui la probabilità cumulata raggiunge o supera} $q$ (e.g. $0.05$, $0.01$).
\begin{equation}
    s_{(k)} = min_{x}\{ x\in \{0,1, 2, \dots, n_{(k)}\}\ |\ \mathbf{F}_{BetaBin}(x;\ n_{(k)},\alpha_{(k)}, \beta_{(k)})\ge q\}
\end{equation}
Trovare analiticamente questo valore risulta assai complicato in quanto $\mathbf{F}_{BetaBin}$ si esprime tramite una \textit{Ipergeometrica Generalizzata} e derivarla è formalmente complesso. Tuttavia possiamo effettuare un'approssimazione. Infatti nel deep learning il numero di neuroni per layer spesso risulta molto grande, sfruttando questa caratteristica, possiamo far riferimento al comportamento asintotico della distribuzione Beta-Binomiale. Per $n_{(k)}$ molto grande la distribuzione Beta-Binomiale tende alla Normale da cui possiamo ottenere una formula analitica \textit{diretta} per $s_{(k)}$
\begin{equation}\label{soglia sk}
    \lceil s_{(k)}\rceil \approx \hat{\mu}_{(k)} + \Phi^{-1}_q\cdot\hat{\sigma}_{(k)} \qquad \text{con } n_{(k)}\rightarrow \infty
\end{equation}
da cui
\begin{align}
    \hat{\mu}_{(k)} &= \mathbb{E}[\hat{L}_{(k)}] = n_{(k)}\dfrac{\hat{\alpha}_{(k)}}{\hat{\alpha}_{(k)}+\hat{\beta}_{(k)}}\\
    \hat{\sigma}_{(k)} &= \sqrt{\mathbb{V}[\hat{L}_{(k)}]} = \sqrt{n_{(k)}\dfrac{\hat{\alpha}_{(k)}[n_{(k)}(1+\hat{\alpha}_{(k)})+\hat{\beta}_{(k)}]}{(\hat{\alpha}_{(k)}+\hat{\beta}_{(k)})(1+\hat{\alpha}_{(k)}+\hat{\beta}_{(k)})}}
\end{align}
$s_{(k)}$ non è più un numero teorico ma bensì \textit{il comportamento minimo osservato in condizioni di successo nel layer-k}.

\subsection{Analisi dell'affidabilità}
Data la definizione di affidabilità [\ref{Affidabilità Layer}], un layer-\textit{k} rappresentato come $L_{(k)}\sim BetaBin(n, \alpha, \beta)$, i parametri $\hat{\alpha}_{(k)}$, $\hat{\beta}_{(k)}$ e la soglia minima $s_{(k)}$; l'\textit{affidabilità} del layer è descritta dalla probabilità:
\begin{equation}
    \mathcal{R}_{(k)}(n_{(k)}, \hat{\alpha}_{(k)}, \hat{\beta}_{(k)}) = \mathbb{P}\{L_{(k)}\ge s_{(k)}\ |\ \hat{\alpha}_{(k)}, \hat{\beta}_{(k)}\} = 1-q
\end{equation}
Stiamo dicendo che nell'analisi della rete imponiamo \textit{a priori} un valore di affidabilità pari a $[(1-q)\cdot 100]\%$ per ogni layer, ciascuno dei quali \textit{garantirà tale affidabilità attivando un numero di neuroni almeno pari ad} $s_{(k)}$. Da qui, l'analisi sull'affidabilità si concentra \textit{a posteriori} nel vedere il numero $s_{(k)}$ di neuroni che ogni layer attiva per garantire l'affidabilità imposta.

Osserviamo la metrica $\rho_{(k)}$ costruita per rappresentare il comportamento del layer
\begin{equation}
    \rho_{(k)} = \bigg[\dfrac{s_{(k)}}{n_{(k)}}\cdot100\bigg]\%
\end{equation}
Questa indica la percentuale minima di neuroni attivi che il layer deve avere per garantire un'affidabilità pari ad $[(1-q)\cdot 100]\%$. Ad esempio, se $\rho_{(k)}=10\%$ con $q=0.05$ allora diremo che il layer-\textit{k} per garantire un affidabilità del $95\%$ deve attivare almeno il $10\%$ dei suoi neuroni.

\textit{Come si può interpretare questa misura?}\\
Un basso valore di $\rho_{(k)}$ indica che al layer basta attivare un piccolo sotto-gruppo di neuroni per garantire di essere statisticamente significativo al $[(1-q)\cdot100]\%$ quando la rete ottiene un successo. Si potrebbe interpretarlo come un segnale di robustezza del layer, in quanto ha molta ridondanza di informazione, d'altra parte un layer con un alto valore di $\rho_{(k)}$ indicherebbe che l'informazione è dispersa tra i neuroni. Tuttavia andrebbero effettuati test per validare tale ipotesi.

\subsection{Intervalli di Confidenza}
Tutto l'approccio descritto fino ad ora ha utilizzato per l'analisi dell'affidabilità, le stime puntuali dei parametri $\alpha$ e $\beta$. Tuttavia, per ottenere una misura più robusta della stima di un parametro è buona norma costruire un \textit{intervallo di confidenza} \footnote{Se $U$ e $V$ sono variabili casuali con distribuzioni di probabilità che dipendono da qualche parametro $\theta$ e $\mathbb{P}[U\le \theta \le V] = 1-\beta$ dove $\beta \in [0, 1]$, allora l'intervallo casuale $(U, V)$ calcolato sul campione osservato, è un intervallo di confidenza al $[(1-\beta)\cdot100]\%$ per $\theta$. I valori estremi dell'intervallo di confidenza si chiamano limiti di confidenza.} in cui tale stima si muoverà con una certa probabilità. Per costruire tale intervallo si possono seguire due strade, una analitica definita da una procedura formale chiamata \textit{Metodo Delta}. La seconda invece, segue la strada computazionale utilizzando \textit{Metodo Bootstrap} stimando gli intervalli da campioni raccolti. In questo lavoro, gli intervalli di confidenza sono derivati usando il \textit{Metodo Delta} \footnote{Il Metodo Delta si compone di 3 fasi: calcolo del gradiente (o matrice Jacobiana) della funzione oggetto dei parametri, costruzione della matrice di covarianza dei parametri con ottenimento della varianza degli stimatori ed infine costruzione dell'intervallo intorno alla stima puntuale dei parametri usando la varianza.}.\\
Tale procedura utilizza il risultato del \textit{Teorema del Limite Centrale} sulla media campionaria $\bar{\mathbf{x}}_{n}$ e la varianza campionaria $S^{2}$, che al crescere di $n$ tendono alla distribuzione \textit{Normale Multivariata}. Il metodo trasferisce questa normalità sulle stime dei parametri $\hat{\alpha}_{(k)},\hat{\beta}_{(k)}$ usando una tecnica di linearizzazione.

\subsubsection{Costruzione IC per $\hat{\alpha}_{(k)}, \hat{\beta}_{(k)}$}
Come prima cosa costruiamo la \textit{funzione vettore} che dobbiamo utilizzare poi nella procedura.
\begin{equation}
    g(m_{1}, m_{2}) = 
    \begin{bmatrix}
        \hat{\alpha}_{(k)}\\
        \hat{\beta}_{(k)}
    \end{bmatrix}
    =
    \begin{bmatrix}
        \dfrac{n_{(k)}m_{1}-m_{2}}{n_{(k)}(\frac{m_{2}}{m_{1}}-{m_{1}} -1)+m_{1}})\\\\
        \dfrac{(n_{(k)}-m_{1})(n_{(k)}-\frac{m_{2}}{m_{1}})}{n_{(k)}(\frac{m_{2}}{m_{1}}-m_{1} -1)+m_{1}}
    \end{bmatrix} 
    =
    \begin{bmatrix}
        \dfrac{m_{1}(nm_{1}-m_{2})}{m_{1}^{2}-n(m_{1}^{2}+m_{1}-m_{2})})\\\\
        \dfrac{(n-m_{1})(nm_{1}-m_{2})}{(1-n)m_{1}^{2} - nm_{1} + nm_{2}}
    \end{bmatrix}
\end{equation}
Dove $m_{1}=\bar{\mathbf{x}}_{n},\ m_{2}=S^{2}$ e l'ultima versione è una riorganizzazione dei termini per non avere frazioni all'interno delle formule.
A partire dalla funzione vettore è necessario costruire la \textit{Matrice Jacobiana} delle derivata parziali.
\begin{equation}
    \mathbf{J}_{2\times2} = \nabla g(m_{1}, m_{2}) = 
    \begin{bmatrix}
        \dfrac{\partial}{\partial m_{1}}\hat{\alpha}_{k} & \dfrac{\partial}{\partial m_{2}}\hat{\alpha}_{k}\\ \\
        \dfrac{\partial}{\partial m_{1}}\hat{\beta}_{k} & \dfrac{\partial}{\partial m_{2}}\hat{\beta}_{k}
    \end{bmatrix}_{2\times2}
\end{equation}
Le derivate parziali da risolvere per $\hat{\alpha}_{(k)}, \hat{\beta}_{(k)}$ sono delle derivate di rapporti, quindi si utilizza la regola di derivazione $\Bigl(\frac{\partial}{\partial x}\frac{f(x)}{g(x)} = \frac{f'(x)g(x)-f(x)g'(x)}{(g(x))^{2}} \Bigr)$. Seguendo questa procedura e fattorizzando le espressioni (per avere una forma quanto più semplice possibile), si ottiene la seguente matrice Jacobiana $\mathbf{J}_{2\times2}$. 
\begin{equation}
    \begin{bmatrix}
        \dfrac{m_{1}^{2}(1-n_{(k)})(m_{2}-n_{(k)}^{2}) - n_{(k)}(n_{(k)}m_{1}-b)^{2}}{(m_{1}^{2}(1-n_{(k)}) - n_{(k)}(m_{1}-m_{2}))^{2}} & 
        \dfrac{m_{1}^{2}(n_{(k)}-1)(m_{1}-n_{(k)})}{(m_{1}^{2}(1-n_{(k)}) - n_{(k)}(m_{1}-m_{2}))^{2}} \\ \\
        \dfrac{m_{1}(1-n_{(k)})(m_{1}-n_{(k)})^{2}}{(m_{1}^{2}(1-n_{(k)}) - n_{(k)}(m_{1}-m_{2}))^{2}} & 
        \dfrac{n_{(k)}(n_{(k)}m_{1}-m_{2})^{2} - m_{2}(1-n_{(k)})(m_{1}-n_{(k)})^{2}}{(m_{1}^{2}(1-n_{(k)}) - n_{(k)}(m_{1}-m_{2}))^{2}}
    \end{bmatrix}_{2\times2}
\end{equation}

Una volta ottenuta $\mathbf{J}_{2\times2}$ è necessario costruire la \textit{Matrice delle Covarianze} $\large\Sigma_{2\times2}$ della media campionaria e varianza campionaria \footnote{Nelle distribuzioni Beta-Binomiali, sono spesso asimmetriche, dunque il termine $\hat{\mu}_{3}$ è diverso da zero. Questo significa che un errore nella stima della media si propaga causando un errore nella stima della varianza. La matrice cattura questa relazione, rendendo l'intervallo di confidenza finale molto più accurato.}. Nel caso dello studio che stiamo portando avanti, la matrice di varianza-covarianza è costruita sui valori osservati durante il \textit{training} della rete. Si costruisce il vettore $\mathbf{x} = \{x_{1}, x_{2}, \dots, x_{N}\}$ del numero di attivazioni \textit{osservate} nel layer-\textit{k} quando la rete restituisce l'output corretto, con $N$ numero totale di predizioni corrette effettuate durante il training. I termini della matrice dipendendo dai \textit{momenti centrati campionari} della distribuzione dei dati $\mathbf{x}$.
\begin{equation}
    \hat{\mu}_{j} = \dfrac{1}{N}\sum_{i=1}^{N}(x_{i}-m_{x})^{j}\qquad\mbox{dove } m_{x} = \dfrac{1}{N}\sum_{i=1}^{N}x_{i} \quad\mbox{(media del campione)}
\end{equation}
La matrice $\large\Sigma_{2\times2}$ si popola utilizzando proprio i momenti centrati campionari normalizzati per la numerosità campionaria $N$.
\begin{equation}
    \Sigma_{2\times2} = \dfrac{1}{N} 
    \begin{bmatrix}
        \hat{\mu}_{2} & \hat{\mu}_{3} \\ \\
        \hat{\mu}_{3} & \hat{\mu}_{4} - (\hat{\mu}_{2})^{2}
    \end{bmatrix}_{2\times2}
\end{equation}
Per costruire l'intervallo di confidenza intorno alle stime puntuali $\hat{\alpha}_{k}, \hat{\beta}_{k}$ è necessario ottenere la \textit{varianza degli stimatori}. Tali valori si ottengono effettuando la moltiplicazione matriciale \textit{sandwich} tra le matrici $\mathbf{J}_{2\times2}$ ed $\large\Sigma_{2\times2}$ e prendendo i valori sulla diagonale principale della matrice risultante.
\begin{equation}\label{matriceSandwich}
     \mathbf{J}_{2\times2}\ \Sigma_{2\times2}\ (\mathbf{J}_{2\times2})^{T} =
     \begin{bmatrix}
         \mathbb{V}(\hat{\alpha}_{(k)}) & \text{Cov}(\hat{\alpha}_{(k)}, \hat{\beta}_{(k)}) \\ \\
         \text{Cov}(\hat{\alpha}_{(k)}, \hat{\beta}_{(k)}) & \mathbb{V}(\hat{\beta}_{(k)})
     \end{bmatrix}_{2\times2}
\end{equation}
la matrice che si ottiene non è altro che la matrice \textit{varianza-covarianza} degli stimatori $\hat{\alpha}_{k}, \hat{\beta}_{k}$. 

Ora possiamo costruire gli intervalli di confidenza asintotici per le stime puntuali facendole variare in modo dipendente dalla loro varianza stimata.
\begin{equation}
    IC_{q}[\alpha_{(k)}]
    = \hat{\alpha}_{(k)}\pm \Phi^{-1}_q\sqrt{\mathbb{V}(\hat{\alpha}_{(k)})} \qquad \quad
    IC_{q}[\beta_{(k)}]
    = \hat{\beta}_{(k)}\pm \Phi^{-1}_q\sqrt{\mathbb{V}(\hat{\beta}_{(k)})}
\end{equation}
dove $\mathbb{V}(\hat{\alpha}_{(k)}) = \Bigl[\mathbf{J}_{2\times2}\ \Sigma_{2\times2}\ (\mathbf{J}_{2\times2})^{T}\Bigr]_{1,1}$ e $\mathbb{V}(\hat{\beta}_{(k)}) = \Bigl[\mathbf{J}_{2\times2}\ \Sigma_{2\times2}\ (\mathbf{J}_{2\times2})^{T}\Bigr]_{2,2}$.

\subsubsection{Costruzione IC per $R_{(k)}$}
Come visto nell'analisi [\ref{soglia sk}], abbiamo studiato l'affidabilità $\mathcal{R}_{(k)}$ di un layer osservando la stima del numero minimo di neuroni che devono essere attivi $s_{(k)}$ per garantire un'affidabilità imposta a priori. Tuttavia, sappiamo che questa dipende anche dalle stime puntuali $\hat{\alpha}_{(k)}, \hat{\beta}_{(k)}$ [\ref{Affidabilità Layer}], quindi è necessario studiare in che intervallo $\mathcal{R}_{(k)}$ si muove al variare di $\hat{\alpha}_{(k)}, \hat{\beta}_{(k)}$.

Come prima cosa otteniamo la \textit{Sensibilità} di $\mathcal{R}_{(k)}$ ossia $\nabla \mathcal{R}_{(k)}(n_{(k)}, \alpha_{(k)}, \beta_{(k)})$ che ci dice come $\mathcal{R}_{(k)}$ si muove al variare dei parametri $\alpha_{(k)}, \beta_{(k)}$.
\begin{equation} \label{vettoreSensibilità}
    \mathbf{R}_{1\times2} = \nabla \mathcal{R}_{(k)}(n_{(k)}, \alpha_{(k)}, \beta_{(k)}) = 
    \begin{bmatrix}
        \dfrac{\partial}{\partial \alpha} \mathcal{R}_{(k)}(n_{(k)}, \alpha_{(k)}, \beta_{(k)}) & \dfrac{\partial}{\partial \beta} \mathcal{R}_{(k)}(n_{(k)}, \alpha_{(k)}, \beta_{(k)})
    \end{bmatrix}
\end{equation}
Consideriamo la formula dell'affidabilità $\mathcal{R}_{(k)}(n_{(k)}, \alpha_{(k)}, \beta_{(k)})$ [\ref{Affidabilità Layer}], sviluppando gli integrali utilizzando la regola $\int_{0}^{1}x^{S}dx = \frac{1}{S+1}$, effettuando i calcoli e riorganizzando i termini si ottiene una nuova forma di $\mathcal{R}_{(k)}$ più facile da maneggiare per costruire $\mathbf{R}_{1\times2}$.
\begin{equation}
    \mathcal{R}_{(k)}(n_{(k)}, \alpha_{(k)}, \beta_{(k)}) = \sum_{j = s_{(k)}}^{n_{(k)}} \binom{n_{(k)}}{j} 
    \dfrac{\frac{1}{A(\alpha_{(k)})} - \frac{1}{B(\alpha_{(k)}, \beta_{(k)})}}{\frac{1}{C(\alpha_{(k)})} - \frac{1}{D(\alpha_{(k)}, \beta_{(k)})}}
\end{equation}
dove $A(\alpha_{(k)}) = \alpha_{(k)} + j,
    B(\alpha_{(k)}, \beta_{(k)}) = \alpha_{(k)} + \beta_{(k)} + n - 1,
    C(\alpha_{(k)}) = \alpha_{(k)} \text{ e }
    D(\alpha_{(k)}, \beta_{(k)}) = \alpha_{(k)} + \beta_{(k)} - 1$.
Utilizzando questa formulazione, nonostante risulti meno complessa della forma originale, per ottenere le derivate parziali per $\alpha_{(k)}, \beta_{(k)}$ più semplici possibili, è necessario riorganizzare l'espressione nella forma $Kf(x)$ in modo tale da avere un termine costante $K$ che non dipendesse dal parametro per cui si sta derivando. In questo modo, sia i calcoli che i risultati finali sono nella forma più semplice possibile.

Partiamo dalla derivata di $\mathcal{R}_{(k)}(n_{(k)}, \alpha_{(k)}, \beta_{(k)})$ rispetto $\alpha_{(k)}$.\\
L'affidabilità espressa in forma $Kf(x)$ è la seguente:
\begin{equation}
    \mathcal{R}_{(k)}(n_{(k)}, \alpha_{(k)}, \beta_{(k)}) = \sum_{j = s_{(k)}}^{n_{(k)}} \binom{n_{(k)}}{j} K_{1}f_{1}(\alpha_{(k)})
\end{equation}
con
\begin{equation}
    K_{1} = \dfrac{\beta_{(k)} + n_{(k)} - j - 1}{\beta_{(k)} - 1} \qquad
    f_{1}(\alpha_{(k)}) = \dfrac{\alpha_{(k)}(\alpha_{(k)}+\beta_{(k)} -1)}{(\alpha_{(k)} + j)(\alpha_{(k)}+\beta_{(k)}+n-1)}
\end{equation}
Calcoliamo la derivata parziale per $\alpha_{(k)}$. Invece di derivare usando la regola classica del rapporto (risulterebbe assai complesso), utilizzo la regola del logaritmo $\frac{\partial}{\partial x}f(x) = f(x)\frac{\partial}{\partial x}\ln(f(x))$, che risulta essere più semplice da derivare.
\begin{equation}
    \dfrac{\partial}{\partial \alpha_{(k)}}\mathcal{R}_{(k)}= \sum_{j = s_{(k)}}^{n_{(k)}} \binom{n_{(k)}}{j} K_{1}f_{1}(\alpha_{(k)})
    \Bigl[ \frac{1}{\alpha_{(k)}} + \frac{1}{\alpha_{(k)} + \beta_{(k)} -1} -\frac{1}{\alpha_{(k)}+j} - \frac{1}{\alpha_{(k)} + \beta_{(k)} + n -1}\Bigr]
\end{equation}

Passiamo alla derivata di $\mathcal{R}_{(k)}(n_{(k)}, \alpha_{(k)}, \beta_{(k)})$ rispetto $\beta_{(k)}$.\\
L'affidabilità espressa in forma $Kf(x)$ è la seguente:
\begin{equation}
    \mathcal{R}_{(k)}(n_{(k)}, \alpha_{(k)}, \beta_{(k)}) = \sum_{j = s_{(k)}}^{n_{(k)}} \binom{n_{(k)}}{j} K_{2}f_{2}(\beta_{(k)})
\end{equation}
Con
\begin{equation}
    K_{2} = \dfrac{\alpha_{(k)}}{\alpha_{(k)} + j} \qquad
    f_{2}(\beta_{(k)}) = \dfrac{(\alpha_{(k)} + \beta_{(k)} - 1)(\beta_{(k)} + n - j - 1)}{(\beta_{(k)}-1)(\alpha_{(k)}+ \beta_{(k)} + n -1)}
\end{equation}
Calcoliamo la derivata parziale per $\beta_{(k)}$, anche qui utilizzo la regola del logaritmo per derivare la funzione.
\begin{equation}
    \dfrac{\partial}{\partial \beta_{(k)}}\mathcal{R}_{(k)}= \sum_{j = s_{(k)}}^{n_{(k)}} \binom{n_{(k)}}{j} K_{2}f_{2}(\beta_{(k)})
    \Bigl[ \frac{1}{\alpha_{(k)} + \beta_{(k)} - 1} + \frac{1}{\beta_{(k)} + n - j - 1} -\frac{1}{\alpha_{(k)}+\beta_{(k)} + n -1} - \frac{1}{\beta_{(k)} - 1}\Bigr]
\end{equation}
Abbiamo ottenuto il vettore \textit{sensibilità} $\mathbf{R}_{1\times2}$ [\ref{vettoreSensibilità}] di $\mathcal{R}_{(k)}$ rispetto $\alpha_{(k)}, \beta_{(k)}$.

Seguendo il metodo Delta, bisogna effettuare la moltiplicazione matriciale \textit{sandwich} tra $\mathbf{R}_{1\times2}$ e la matrice di covarianza delle stime dei parametri $\alpha_{(k)}, \beta_{(k)}$. Ma questa matrice già l'abbiamo ottenuta durante la costruzione dell'intervallo di confidenza dei parametri quando abbiamo effettuato la moltiplicazione matriciale $\mathbf{J}_{2\times2}\ \Sigma_{2\times2}\ (\mathbf{J}_{2\times2})^{T}$ [\ref{matriceSandwich}]. Indicando con $\hat{\Sigma}_{2\times2}$ tale matrice, possiamo ottenere la \textit{varianza} di $\mathcal{R}_{(k)}(n_{(k)}, \alpha_{(k)}, \beta_{(k)})$.
\begin{equation}
    \mathbb{V}[\mathcal{R}_{(k)}(n_{(k)}, \alpha_{(k)}, \beta_{(k)})] \approx (\mathbf{R}_{1\times2})^{T}\ \hat{\Sigma}_{2\times2} \ (\mathbf{R}_{1\times2})  
\end{equation}
Il risultato è uno scalare che rappresenta la varianza dell'affidabilità $\mathcal{R}_{(k)}(n_{(k)}, \alpha_{(k)}, \beta_{(k)})$ del layer-\textit{k}. Ottenuto questo valore è possibile costruire l'\textit{intervallo di confidenza} dell'affidabilità intorno alla sua stima nei punti $\hat{\alpha}_{(k)},\hat{\beta}_{(k)}$.
\begin{equation}
    IC_{q}[\mathcal{R}_{(k)}(n_{(k)}, \alpha_{(k)}, \beta_{(k)})] = \mathcal{R}_{(k)}(n_{(k)}, \hat{\alpha}_{(k)}, \hat{\beta}_{(k)}) \pm \Phi^{-1}_{q}\sqrt{\mathbb{V}[\mathcal{R}_{(k)}(n_{(k)}, \alpha_{(k)}, \beta_{(k)})]}
\end{equation}
Questo intervallo di confidenza ci permette di quantificare l'\textit{incertezza} associata alla stima puntuale $\mathcal{R}_{(k)}$, così da poter effettuare una valutazione critica delle reti neurali come si fa in contesti \textit{safety-critical} per i sistemi classici.
\begin{enumerate}
    \item Identificazione del \textit{Lower Confidence Bound}. \\
    Nell'analisi safety-critical dei sistemi si è interessati al \textit{caso peggiore}. In questo caso, l'estremo inferiore dell'intervallo rappresenta la \textit{minima affidabilità garantita} con un certo livello di confidenza specificato (e.g. $95\%$). Quindi, una volta ottenuto l'intervallo, il limite inferiore ci permette di affermare che: \textit{Il layer-k, con una soglia di attivazione minima di neuroni pari ad $s_{(k)}$, garantisce un'affidabilità almeno pari all'estremo inferiore dell'intervallo con una confidenza del $[(1-q)\cdot100]\%$}\footnote{Ipotizziamo che dopo il training, seguendo la procedura descritta fino ad ora si stimano $\hat{\alpha}_{(k)}, \hat{\beta}_{(k)}$, si ottiene la metrica $\rho_{(k)} = [s_{(k)}/{n_{(k)}}\cdot100]\% = 40\%$ per il layer-\textit{k}. Costruendo l'intervallo di confindenza intorno al valore dell'affidabilità $R_{(k)}$ in $\hat{\alpha}_{(k)}, \hat{\beta}_{(k)}$ con una confidenza del $95\%$, si osserva un limite inferiore pari a circa il $75\%$. Allora diremo che, il layer-\textit{k} avrà un'affidabilità almeno del $75\%$, nel $95\%$ delle volte, attivando circa il $40\%$ dei suoi neuroni.}.
    
    \item Metrica di \textit{robustezza dell'addestramento}.\footnote{Lo studio di come l'intervallo si comporta asintoticamente e in funzione di quali valori verrà affrontato nei prossimi lavori.}\\
    L'ampiezza dell'intervallo potrebbe esser vista anche come una misura legata alla qualità del training per uno specifico layer.
    \begin{itemize}
        \item Intervallo stretto implica \textit{meno} variabilità e quindi una convergenza \textit{stabile};
        \item Intervallo largo indica che indipendentemente dall'accuratezza, il comportamento interno al layer è \textit{volatile}.
    \end{itemize}
\end{enumerate}

\section{Conclusioni}
Lo studio è un tentativo di introdurre un metodo per analizzare le reti neurali, proponendo un passaggio dalla semplice valutazione performance-based ad una valutazione \textit{reliability-based}. Modellando i layer come sistemi stocastici \textit{k}-out-of-\textit{n} rappresentati da distribuzioni Beta-Binomiali, è possibile quantificare la \textit{partecipazione dei neuroni} e definire soglie di attivazione data-driven come $s_{(k)}$ e metriche operative come $\rho_{(k)}$.

% Tuttavia, il contributo più significativo di questo approccio, a mio parere, risiede nel suo utilizzo per l'Explainable AI (XAI). Siccome le tecniche tradizionali di XAI si concentrano nell'analizzare \textit{perché} una rete ha preso una determinata decisione, affiancarci questo approccio significa rispondere anche alla domanda: \textit{Con quanta affidabilità interna la rete sta operando?}
 
% Poiché questo lavoro si è focalizzato sulla definizione e la dimostrazione teorica della metodologia, i passi successivi che mi sono posto sono
% \begin{itemize}
%     \item Ampliare lo studio teorico, studiando le funzioni che sono state costruite per l'affidabilità e gli intervalli di confidenza, per dedurne altre caratteristiche strutturali;
%     \item Sviluppo di un framework applicativo capace di \textit{automatizzare l'inferenza statistica}, \textit{monitoraggio in runtime} dell'affidabilità dei layer e la loro \textit{validazione}.
% \end{itemize}

\end{document}
